\documentclass[aspectratio=169,10pt]{beamer}
\usetheme[numbering=fraction,progressbar=frametitle]{metropolis}
\usepackage{fontspec}
\setsansfont{Latin Modern Sans}
\usepackage{amsmath,amssymb,mathtools,graphicx,booktabs}
\usepackage{polyglossia}
\setdefaultlanguage{french}
\setbeamertemplate{navigation symbols}{}
\setbeamerfont{frametitle}{size=\large}
\setbeamerfont{normal text}{size=\normalsize}
\setlength{\parskip}{.45em}
\hypersetup{pdfauthor={},pdfsubject={Lecture slides},pdfkeywords={},pdftitle={Vecteurs du plan}}
\providecommand{\N}{\mathbb{N}}
\providecommand{\Z}{\mathbb{Z}}
\providecommand{\D}{\mathbb{D}}
\providecommand{\Q}{\mathbb{Q}}
\providecommand{\R}{\mathbb{R}}
\providecommand{\C}{\mathbb{C}}
\providecommand{\vect}[1]{\overrightarrow{#1}}
\providecommand{\vi}{\vec{\imath}}
\providecommand{\vj}{\vec{\jmath}}
\providecommand{\vk}{\vec{k}}
\providecommand{\norm}[1]{\left\lVert #1\right\rVert}
\providecommand{\vc}[2]{\begin{pmatrix}#1\\#2\end{pmatrix}}
\providecommand{\deter}[2]{\det\left(#1,#2\right)}
\providecommand{\abs}[1]{\left\lvert #1\right\rvert}
\providecommand{\dd}{\,\mathrm{d}}
\providecommand{\e}{\mathrm{e}}
\providecommand{\Card}{\operatorname{Card}}
\providecommand{\E}{\operatorname{E}}
\providecommand{\V}{\operatorname{V}}
\providecommand{\degre}{^{\circ}}
\providecommand{\un}[1]{\,\mathrm{#1}}
\providecommand{\sci}[2]{#1\times10^{#2}}
\providecommand{\pgcd}{\operatorname{PGCD}}
\providecommand{\ppcm}{\operatorname{PPCM}}
\title{Vecteurs du plan}
\subtitle{Cours et exemples}
\date{}
\author{}
\begin{document}
\begin{frame}[plain]
\vfill
{\usebeamerfont{title}\inserttitle\par}
\medskip
{\usebeamerfont{subtitle}\insertsubtitle\par}
\medskip
{\usebeamercolor[fg]{progress bar}\rule{\textwidth}{.5pt}}
\vfill
\end{frame}
\begin{frame}{}
Sur un plan quadrillé, plusieurs chemins conduisent du point \(A\) au point \(B\).
Coder chacun par ses déplacements horizontal et vertical. Identifier ce qui
reste invariant malgré le chemin choisi, puis proposer une règle d'addition de
deux déplacements successifs.
\end{frame}
\begin{frame}{Objectifs}
\begin{itemize}

  \item Connaître les caractéristiques d'un vecteur (direction, sens, norme).
  
\item Maîtriser l'addition vectorielle et la multiplication par un scalaire.
  
\item Calculer et utiliser le produit scalaire.
  
\item Démontrer colinéarité et orthogonalité.

\end{itemize}
\end{frame}
\begin{frame}{Généralités}
Un \textbf{vecteur non nul} \(\overrightarrow{AB}\) est défini par sa \textbf{direction},
son \textbf{sens} et sa \textbf{norme} \(\left\|\overrightarrow{AB}\right\| = AB\).
Si \(A=B\), \(\overrightarrow{AB}=\vec0\) : sa norme est nulle et il n'a ni direction ni sens.
Deux vecteurs non nuls sont \textbf{égaux} s'ils ont même direction, sens et norme.
\end{frame}
\begin{frame}{Généralités}
\(I\) est le milieu de \([AB]\) si et seulement si
\(\overrightarrow{IA} + \overrightarrow{IB} = \overrightarrow{0}\).
\end{frame}
\begin{frame}{Produit scalaire}
Soient \(\vec{u}=(x_1,y_1)\) et \(\vec{v}=(x_2,y_2)\) :
\[\vec{u} \cdot \vec{v} = x_1 x_2 + y_1 y_2
  .\]
Si les deux vecteurs sont non nuls, on a aussi
\(\vec u\cdot\vec v=\left\|\vec u\right\|\,\left\|\vec v\right\|\cos\theta\),
où \(\theta\) est leur angle. Si l'un est nul, le produit scalaire vaut \(0\).
\end{frame}
\begin{frame}{Produit scalaire}
\begin{itemize}

  \item \(\vec{u} \perp \vec{v} \Leftrightarrow \vec{u}\cdot\vec{v}=0\).
  
\item \(\vec{u} \parallel \vec{v} \Leftrightarrow x_1 y_2 - x_2 y_1 = 0\).

\end{itemize}
\end{frame}
\begin{frame}{Opérations sur les vecteurs}
Soient \(\vec{u}=(x_1,y_1)\), \(\vec{v}=(x_2,y_2)\) et \(\lambda\in\mathbb{R}\).
\[\vec{u}+\vec{v} = (x_1+x_2,\, y_1+y_2),\qquad
  \lambda\vec{u} = (\lambda x_1,\, \lambda y_1).\]
Le vecteur \(-\vec{u} = (-x_1,-y_1)\) est l'\textbf{opposé} de \(\vec{u}\).
\end{frame}
\begin{frame}{Opérations sur les vecteurs}
\begin{itemize}

  \item \textbf{Commutativité :} \(\vec{u}+\vec{v}=\vec{v}+\vec{u}\).
  
\item \textbf{Associativité :} \((\vec{u}+\vec{v})+\vec{w}=\vec{u}+(\vec{v}+\vec{w})\).
  
\item \textbf{Élément neutre :} \(\vec{u}+\vec{0}=\vec{u}\).
  
\item \textbf{Relation de Chasles :} \(\overrightarrow{AB}+\overrightarrow{BC}=\overrightarrow{AC}\).

\end{itemize}
\end{frame}
\begin{frame}{Opérations sur les vecteurs}
Si \(A(x_A,y_A)\) et \(B(x_B,y_B)\), alors :
\[\overrightarrow{AB} = (x_B-x_A,\; y_B-y_A),\qquad
  \left\|\overrightarrow{AB}\right\| = \sqrt{(x_B-x_A)^2+(y_B-y_A)^2}.\]
\end{frame}
\begin{frame}{Opérations sur les vecteurs}
\(\vec{u}=(x_1,y_1)\) et \(\vec{v}=(x_2,y_2)\) sont \textbf{colinéaires}
si et seulement si :
\[x_1 y_2 - x_2 y_1 = 0.\]
Deux vecteurs non nuls colinéaires ont la même direction.
Le vecteur nul est colinéaire à tout vecteur.
Cette expression s'appelle le \textbf{déterminant} du couple \((\vec{u},\vec{v})\).
\end{frame}
\begin{frame}{Opérations sur les vecteurs}
Le milieu \(I\) de \([AB]\) a pour coordonnées :
\(I\!\left(\dfrac{x_A+x_B}{2},\; \dfrac{y_A+y_B}{2}\right).\)

Le centre de gravité \(G\) du triangle \(ABC\) vérifie :
\(\overrightarrow{GA}+\overrightarrow{GB}+\overrightarrow{GC}=\vec{0}\), et ses coordonnées sont
\(G\!\left(\dfrac{x_A+x_B+x_C}{3},\; \dfrac{y_A+y_B+y_C}{3}\right).\)
\end{frame}
\begin{frame}{Repère et coordonnées}
Un \textbf{repère orthonormé} \((O, \vec{\imath}, \vec{\jmath})\) est formé d'un
point \(O\) (origine) et de deux vecteurs perpendiculaires de norme \(1\).
Tout vecteur \(\vec{u}\) s'écrit \(\vec{u} = x\vec{\imath} + y\vec{\jmath}\) ;
\((x,y)\) sont les \textbf{coordonnées} de \(\vec{u}\).
\end{frame}
\begin{frame}{Repère et coordonnées}
\(A\), \(B\), \(C\) sont alignés ssi \(\overrightarrow{AB}\) et \(\overrightarrow{AC}\) sont colinéaires, c'est-à-dire :
\[\det(\overrightarrow{AB}, \overrightarrow{AC})
  = (x_B-x_A)(y_C-y_A) - (x_C-x_A)(y_B-y_A) = 0.\]
\end{frame}
\begin{frame}{Repère et coordonnées}
\(ABCD\) est un parallélogramme ssi \(\overrightarrow{AB} = \overrightarrow{DC}\)
(ou, de manière équivalente, si les diagonales se coupent en leur milieu).
\end{frame}
\end{document}
